\documentclass[conference]{IEEEtran}
\IEEEoverridecommandlockouts

\usepackage{placeins}
\usepackage{cite}
\usepackage{amsmath,amssymb,amsfonts}
\usepackage{algorithmic}
\usepackage{graphicx}
\usepackage{textcomp}
\usepackage{xcolor}

\begin{document}

\title{Federated Learning-Enabled Secure Quantum Key Distribution: A Simulation-Based Study}

\author{

\IEEEauthorblockN{Raman Sharma}
\IEEEauthorblockA{
Department of Electronics and Communication Engineering\\
National Institute of Technology Kurukshetra\\
Kurukshetra, India\\
ramansharma05384@gmail.com}

\and

\IEEEauthorblockN{Pankaj Verma}
\IEEEauthorblockA{
Department of Electronics and Communication Engineering\\
National Institute of Technology Kurukshetra\\
Kurukshetra, India\\
pankaj@nitkkr.ac.in}

}

\maketitle

\begin{abstract}

Quantum Key Distribution (QKD) is the one that offers a secure exchange mechanism of cryptographic keys executed on the basis of the principles of quantum mechanics. Practical QKD systems are however challenged by hardware imperfection, photon loss, and channel noise are some of the challenges and variation of Quantum Bit Error Rate (QBER). There are factors that diminish the efficiency of secret key generation in real-world implementations.In this paper, a simulation-based presentation is made model combining a QKD model based on BB84 with Federated Learning (FL). The suggested system produces quantum communication data in various channel conditions such as noise, attenuation, and transmission distance. Distributed nodes train local models of machine learning without losing the privacy.The locally trained models are of their raw measurement data generated using the federated averaging algorithm to aggregate a universal tranquilizer that can forecast the values of QBER and identify them possible eavesdropping attacks. The outcomes of simulation prove that the proposed federated learning system attains prediction performance in the same manner as centralized learning methods and at the same time ensuring the privacy of data and minimizing communication overhead.

\end{abstract}


Keywords— Quantum Key Distribution (QKD), BB84 Protocol, Federated Learning (FL), Quantum Bit Error Rate (QBER), Distributed Machine Learning, Privacy Preservation, Eavesdropping Detection, Quantum Communication, Federated Averaging (FedAvg), Secure Key Generation.


\section{Introduction}

In the contemporary communication systems, secure communication is a requirement. The traditional cryptographic approaches are based on computational complexity as their security, which can be compromised due to the development of quantum computing devices. The communication approach of quantum communication has a promising future where quantum mechanics is used to realize theoretically secure communication.

Quantum Key Distribution (QKD) facilitates sharing of encryption keys between two parties with the help of quantum bits in quantum channel. The BB84 protocol is one of the most popular protocols that encodes the information by using quantum states of photons. Any form of eavesdropping presents some detectable interference in the communication channel.There are however, various constraints affecting practical QKD systems which include noise on the channel, loss of photons and imperfections of the hardware. Such problems lead to fluctuations in the Quantum Bit Error Rate (QBER) of which the reliability of key generation directly depends.

The concept of machine learning has been recently pursued to forecast the behavior of the channel and enhance the reliability of communication. Federated Learning (FL) is a technology in distributed machine learning in which more than one node trains a global model together without exchanging raw data.

The main contributions of this work are summarized as follows:

\begin{itemize}

\item Design of a BB84-based QKD communication model.

\item Integration of federated learning for distributed QBER prediction.

\item Privacy-preserving model aggregation using federated learning.

\item Performance evaluation using MATLAB-based simulations.

\end{itemize}

\section{Background and Related Work}

\subsection{Federated Learning Frameworks}

Federated Learning (FL) is a distributed machine learning approach, in which a number of clients collaboratively trains a global model using their respective local datasets that remain confidential. As opposed to sending raw data to a central server, every client learns a local model and sends model parameters, which enhances privacy and minimizes the communication overhead. The architecture removes the central server in decentralized federated learning and enables the nodes to share model parameters directly with the other devices, which enhances the scalability and resilience of distributed networks. \cite{b8}.


Federated learning is semi-decentralized, which implies that the communication between devices and centralized coordination is merged. Local models are then summed up with the surrounding nodes and then sent to a central node which minimizes communication overhead. \cite{b9}.

The article suggests a hybrid wireless federated learning system, which will include local training on privacy sensitive devices and centralized training on resource constrained devices that are not privacy sensitive. The hybrid design can enhance the training efficiency and lower latency through optimization of the device training mode, training parameters and communication resources. In the case of federated learning in Quantum Key Distribution (QKD), this type of framework will provide a compromise solution that secures quantum data privacy through local quantum training and facilitates centralized quantum processing where needed, to support the efficient, secure, and scalable quantum federated learning.  \cite{b11}

The paper studies optimizing federated learning (FL) convergence over wireless networks with limited resources. It proposes a probabilistic user selection that prioritizes users whose local models strongly impact the global model and jointly allocates wireless uplink resources. Additionally, artificial neural networks estimate local models of non-selected users to enhance global model accuracy. This approach significantly reduces FL convergence time and improves training accuracy compared to standard methods. \cite{b12}



\subsection{Privacy and Security in Federated Learning}

Federated learning improves privacy since models are not trained on the raw data. Nonetheless, information leakage can still be done through model updates sent during training. To counter this, privacy-sensitive methods like Private Aggregation of Teacher Ensembles, Homomorphic Encryption and Secure Multi-Party Computer are introduced in conjunction with federated learning. These approaches are used to ensure the model updates and aggregation, balance between privacy and utility, and are generally applied in highly sensitive areas such as healthcare, finance, and IoT to maintain security against malicious attacks and retain data confidentiality (Eman Shalabi, 2025; Kosaraju, 2024; Prudhivi Anuradha et al., 2025).  \cite{b13} 


\subsection {Quantum Federated Learning Approaches}

The emergence of new studies discusses the possibilities of combining federated learning with quantum technologies. It is possible to combine federated learning with quantum communication systems, like QKD to have intelligent, adaptive, and secure communication networks. \cite{b10}.

One of the recent studies explored Quantum Federated Learning (QFL) as a distributed system of training Quantum Machine Learning (QML) models with data privacy. The authors suggested the idea of Federated Quantum Neural Network (FedQNN), which integrates the features of quantum machine learning and classical principles of federated learning. The framework allows joint model training in distributed nodes without exchanging the raw data, thus enhancing privacy and security. The results of experiment on several datasets of healthcare and genomics showed that the proposed model has more than 86 percent classification accuracy and can be successfully applied in a number of quantum machine learning problems. This article underscores the opportunities of quantum computing and federated learning integration as a way of providing safe and collaborative applications of data.
 \cite{b14}


Recently, there was a suggested work on Federated Quantum Learning (FQL) framework, which allows several different quantum nodes to jointly optimize variational quantum circuits without raw data exchange. The framework combines quantum key distribution and secure gradient aggregation protocols to ensure privacy and efficiency in communication. Experimental outcomes on quantum simulators indicated that the given approach could have the same performance as the centralized quantum learning and imposed less than 50 percent communication cost than classical federated learning. This method illustrates how federated quantum learning can be used to implement privacy-sensitive federated quantum learning problems like healthcare, finance and national security.  \cite{b15}


\subsection{Quantum Key Distribution and BB84 Studies}

One of the studies examined the performance of the BB84 protocol as an effective way of securely distributing quantum keys under the circumstance of the possible presence of eavesdroppers. The protocol employs the basic quantum principles like, superposition, entanglement and the no-cloning principle to facilitate safe key exchange between communicating parties. It was implemented and analyzed with the help of Qiskit in which an ideal quantum communication channel was provided. It was experimentally proven that an eavesdropper attempting to measure the qubits with an inappropriate basis causes detectable perturbations in the quantum states and as a consequence, the concerned parties can detect the attempts by the eavesdropper to intercept the communications and ensure the privacy of the communications.  \cite{b16}

\subsection{Federated vs Centralized Learning Performance}

The effect of data quality on centralized and federated learning models was studied through the analysis of datasets that had lower data accuracy and data completeness. The results of the experiments indicated that federated learning had better validation performance than centralized learning, especially when the low-quality data is uniformly distributed among the clients. The results show that federated learning systems tend to tolerate low data quality better than centralized systems, which reflects its suitability to distributed and privacy-preserving machine learning systems. \cite{b17}





\section{System Model}

The suggested system will combine Quantum Key Distribution (QKD) with the BB84 protocol with Federated Learning (FL) to facilitate intelligent and privacy-aware prediction of Quantum Bit Error Rate (QBER). The general design of the proposed architecture is shown in Fig.1.
\begin{figure}[htbp]
    \centering
    \includegraphics[width=\linewidth]{federated system model.png}
    \caption{Proposed BB84-based quantum communication framework integrated with federated learning.}
    \label{fig.1:system_model}
\end{figure}

As illustrated in Fig. 1, several distributed QKD nodes are involved with different channel conditions and they jointly train a global prediction model without any exchange of raw communication information. The overall communication and learning process comprises of the following stages: 

\subsection*{1) Random Bit Generation}

Random binary bits are created at the transmitter side (Alice). These fragments make up raw key material which will be sent via the quantum channel.  

\subsection*{2) BB84 Quantum State Encoding}

The bits that are generated are encrypted by the BB84 protocol. Randomly each bit is assigned to one of 2 polarization bases (rectilinear or diagonal). Such manner of selecting the basis in a random fashion provides security since, any attempt at interception will cause detectable perturbations in the quantum states sent. 

\subsection*{3) Transmission Through Quantum Channel}

The photons are sent along a quantum channel in the form of harmonies. When transmitting, the signal can be impaired by noise in the channels, attenuation, loss of photons, environment disruption, and possible attack of eavesdropping. Such impairments have a direct effect on the performance of the systems. 

\subsection*{4) Measurement at Receiver}

Measuring of the incoming photons is done at the receiver side (Bob) where random selection of bases is used. Once transmitted, Alice and Bob are connected through a classical channel where a basis reconciliation is done to ignore measurements that are not consistent, and only the valid bits of the key are kept. 

\subsection*{5) QBER Calculation}

The Quantum Bit Error Rate (QBER) is calculated by comparing some of the transmitted and the received bits. QBER is the ratio of incorrect bits transferred to the number of bits transferred and is an important performance and security measure. A large QBER can be an indication of too much noise or even eavesdropping.

\subsection*{6) Local Machine Learning Model Training}

Local channel statistics such as the noise level, attenuation, distance of transmission and the QBER measured by each distributed node. Based on this locally obtained dataset a machine learning model is trained to forecast QBER at varying channel conditions. Notably, data in raw form is stored on the local level to guarantee privacy. 

\subsection*{7) Federated Model Aggregation}

Each node does not communicate the raw communication data, instead it only transfers its trained model parameters to a central federated server. These local model update examples are collectively summed by the server through the Federated Averaging (FedAvg) algorithm to create a global QBER prediction model. The revised international model is then re-distributed to all of the nodes to be re-trained in the next round of communication. The model is trained in local and global steps whereby through local training and global aggregation, the model will converge to a correct predictor of QBER and still retain the data confidentiality. The framework of this distributed learning is better at scale, more privacy-enhancing and smarter at monitoring the performance of the QKD system. 

\section{Simulation Setup}

The suggested federated learning-based QKD system is simulated and tested in MATLAB. The simulation aims to evaluate the efficiency of distributed predicting QBER at different conditions of the quantum channel considering that the data privacy is maintained among a number of nodes. 

\subsection*{A. Distributed QKD Node Configuration}
The simulation setup is that of a set of distributed QKD nodes each of which models an independent quantum communication channel under varying channel conditions. Each node is a full-scale simulation of a key distribution that is based on BB84, such as random bit generation, quantum state encoding, transmission over a noisy quantum channel, receiver measurement, and estimation of QBER. 
In order to simulate real world communication conditions, the channel parameters of the different nodes are set to be different including noise variance, attenuation coefficient, and probability of transmitting photons. This guarantees the absence of identical data distributions among nodes which is a real-world distributed network setup. 

\subsection*{B. Quantum Channel Modeling}

The quantum channel is modeled to incorporate the following impairments:

\begin{itemize}
\item Additive channel noise affecting photon polarization states
\item Photon loss due to attenuation over transmission distance
\item Random bit flips to simulate environmental disturbances
\end{itemize}

The Quantum Bit Error Rate (QBER) is calculated at each node as:

\[
QBER = \frac{N_{error}}{N_{total}}
\]

where $N_{error}$ represents the number of incorrectly received bits and $N_{total}$ denotes the total number of transmitted bits after basis reconciliation.

By varying noise levels and transmission probabilities, different QBER distributions are generated across nodes.

\subsection*{C. Local Machine Learning Model Training}

The distributed nodes are trained separately whereby a local prediction model is trained on locally acquired QKD data. This model is aimed at learning how the conditions of quantum channels are related to the resulting Quantum Bit Error Rate (QBER). Privacy is maintained during the training process since raw quantum communication data is not exchanged between nodes still is stored there.
In order to model the nonlinear behavior of QBER under different channel conditions, a feedforward Deep Neural Network (DNN) is deployed at every of the fed clients. 

\subsubsection{ Input Features}
The input layer of the neural network consists of physical parameters derived from the simulated BB84-based QKD channel, including:

\begin{itemize}
    \item Channel noise level
    \item Photon transmission probability
    \item Transmission distance (modeled attenuation)
    \item Detector efficiency
    \item Dark count rate
\end{itemize}

These parameters collectively describe the state of the quantum communication link and directly influence the observed QBER variations.

\subsubsection{ Network Architecture}

The proposed DNN contains two fully connected hidden layers:

\begin{itemize}
    \item First hidden layer: 64 neurons
    \item Second hidden layer: 32 neurons
\end{itemize}
The Rectified Linear Unit (ReLU) activation function is applied to both hidden layers to introduce nonlinearity and enable the model to learn complex feature interactions.

To enhance training stability and improve convergence across heterogeneous client datasets, batch normalization is incorporated after each hidden layer.

\subsubsection{ Output Layer}

The output neuron is a single unit, which produces a continuous-valued prediction of the estimated QBER. Since the formulation of QBER prediction is a regression task, no activation used at the last layer. 

 

\subsubsection{ Training Objective and Optimization}

Its model is trained through supervised learning. The optimization problem is to obtain the minimum of the Mean Absolute Error (MAE) which is the mean value of the difference between the predicted and actual values of QBER. The reason behind the choice of MAE is that it is resistant to noise variations and is appropriate to be used in regression-based performance estimation. 

The Adam optimization algorithm is used to perform parameter updates, and it offers adaptive learning rates and the robust convergence behavior to distributed training. 
\subsubsection{ Federated Learning Integration}

During each communication round, every client trains the DNN locally for a predefined number of epochs. After local training, only the updated model weights are transmitted to the central aggregation server. The global model parameters are computed using the Federated Averaging (FedAvg) algorithm and redistributed to all clients for the next training iteration.

At every communication round, all the clients train the DNN locally on a fixed amount of epochs. Local training is performed and then the new model weights are sent to the central aggregation server. The world model parameters are calculated with Federated Averaging (FedAvg) algorithm and sent back to all clients, to use in the next training step.
It allows QBER prediction in a privacy-preserving manner and allows the correct prediction over distributed QKD nodes without having to collect data centrally. 


\subsection*{D. Federated Learning Procedure}

The trained model parameters (weights and biases) are then only sent to a central federated server by the local training. The federated averaging (FedAvg) algorithm is used by the server to do global model aggregation:

\[
w^{(t+1)} = \sum_{k=1}^{K} \frac{n_k}{n} w_k^{(t)}
\]

where $w_k^{(t)}$ represents the model parameters from node $k$, $n_k$ is the number of local samples at node $k$, and $n$ is the total number of samples across all nodes.

These parameters are varied systematically in order to test convergence behavior, accuracy of prediction as well as efficiency of communication. 

\subsection*{E. Simulation Parameters}

The key simulation parameters used in the experiments include:

\begin{itemize}
\item Number of distributed nodes participating in federated training
\item Channel noise level range
\item Photon transmission probability
\item Number of federated communication rounds
\item Number of local training epochs per round
\end{itemize}

These parameters are systematically varied to evaluate convergence behavior, prediction accuracy, and communication efficiency.

\subsection*{F. Evaluation Metrics}

The performance of the proposed framework is evaluated using:

\begin{itemize}
\item Mean Squared Error (MSE) for QBER prediction
\item Prediction error distribution
\item Convergence rate across federated rounds
\item Comparison with centralized learning baseline
\end{itemize}

In this simulation framework, the proposed system proves that federated learning can be effectively used to predict QBER with a variety of channel conditions and has a high reduction in raw data transmission as well as the local privacy is preserved. 


\section{Results and Discussion}

The effectiveness of the proposed federated learning-based framework in distributed QKD systems to predict QBER is justified by the simulation results. The main aim of the evaluation is to evaluate the accuracy of prediction, convergence behavior, and performance comparison with the centralized learning and the preservation of data privacy. 

The experimental outcomes prove that the federated learning structure is accurate in depicting the interaction between the impairments of quantum channels and the consequent Quantum Bit Error Rate (QBER). The global model also converges steadily and results in high prediction accuracy, even though the distributed and non-identically distributed (non-IID) data is distributed across the nodes. 

 

\subsection*{A. QBER Prediction Performance}

The results of QBER comparison indicate that there is obvious co-relation between noise levels in the channels, transmission probability of the photons, and variations in error rates. As it would be anticipated, channel noise and attenuation increases result in a corresponding rise in QBER. The federated model has managed to learn this nonlinear relationship and give predictions that are in close correspondence with real measured QBER values. 

The fact that the prediction deviation is low shows that the distributed training does not significantly harm performance as opposed to centralized learning. 

 
\subsection*{B. Prediction Error Analysis}

The results of the error distribution indicate that the concentration of prediction errors is around the value zero. This limited distribution validates the fact that the trained global model is generalizable over a variety of distributed nodes with heterogeneous channel conditions. 

The small value of Mean Squared Error (MSE) also shows that the regression was stable and resistant to the variation in noise. 

\subsection*{C. Federated Learning Convergence}

The convergence analysis of the federated communication rounds shows that the training loss is gradually decreasing. First, the model suffers greater loss because of randomization of initialisation and local datasets that are not similar. Nonetheless, the global model gradually becomes stabilized with the repetition of the parameter aggregation by the Federated Averaging algorithm. 

The convergence curve is smooth, which is an affirmation that the distributed knowledge is effectively combined by the aggregation strategy without experiencing the model divergence. 

\subsection*{D. Comparison with Centralized Learning}

 The development of the content and the experimental processes did not encompass the aspects of learning because the centralized way of learning is not detailed.
 
 A comparative study of the federated and centralized methods of learning shows that federated learning has an almost similar prediction accuracy as compared to centralized learning method. Though the centralized learning has a direct access to all training data, the federated framework has a similar performance, but it greatly lowers raw data transmission.
 
 This puts forth the main strength of the suggested strategy: privacy, scalability without obstructing predictive ability. 
\subsection*{E. Classification and Security Implications}

Ending up with the confusion matrix results, the true positive and the true negative rates are high in terms of identifying abnormal channel conditions. This confirms that the suggested framework can effectively determine normal fluctuations in channels as opposed to the situation of eavesdropping.

Proper prediction of QBER allows the QKD systems to be monitored proactively and adaptive security in the real world quantum communication network.

In general, the findings validate that federated learning with support of the BB84-based QKD can offer an intelligent, scalable, and privacy-sensitive solution in the case of a distributed quantum communication setting. 

\subsection*{F. Figures and Table}

%-----------------------------------------
\subsection{QBER Comparison Analysis}

\noindent\textbf{Figure 1: QBER variation under different channel conditions}

\begin{figure}[htbp]
\centering
\includegraphics[width=0.9\columnwidth]{Figure1_QBER_Prediction.png}
\label{fig:qbercomparison}
\end{figure}

\noindent
Figure 1 shows the dependence of QBER on various quantum channel conditions. The findings indicate that the more the channel noise and attenuation, the greater is the QBER. This is because the federated learning model is able to capture such variations and come up with trustworthy QBER predictions. This helps to detect the behavior of abnormal channels and possible eavesdropping of the QKD system. 
\FloatBarrier

%-----------------------------------------
\subsection{QBER Prediction Error Distribution}

\noindent\textbf{Figure 2: Prediction error distribution of the federated learning model}

\begin{figure}[htbp]
\centering
\includegraphics[width=0.9\columnwidth]{Figure2_Error_Distribution.png}
\label{fig:qbererror}
\end{figure}

\noindent
In Figure 2, the distribution of errors in prediction made through the trained federated learning model is presented. The error of prediction is very small around zero, which means that there is high prediction accuracy. The small error distribution verifies great generalization ability between dispersed nodes. 

\FloatBarrier

%-----------------------------------------
\subsection{Federated Learning Convergence}

\noindent\textbf{Figure 3: Federated learning convergence across communication rounds}

\begin{figure}[htbp]
\centering
\includegraphics[width=0.9\columnwidth]{Figure3_Convergence.png}
\label{fig:flconvergence}
\end{figure}

\noindent
Figure 3 demonstrates the convergence of the federated learning process in several communication rounds. The rounds of training loss reduce slowly, which proves that the learning process is stable and the global model is developed continuously. 

\FloatBarrier

%-----------------------------------------
\subsection{Federated Learning vs Centralized Learning}

\noindent\textbf{Figure 4: Performance comparison between federated learning and centralized learning}

\begin{figure}[htbp]
\centering
\includegraphics[width=0.9\columnwidth]{Figure4_FL_vs_Centralized.png}
\label{fig:comparison}
\end{figure}

\noindent
The performance of federated learning and centralized learning methods is compared in Figure\ 4. The findings suggest that federated learning can achieve similar accuracy in prediction and much less raw data has to be transmitted, thus improving privacy. 

\FloatBarrier

%-----------------------------------------
\subsection{Classification Performance}

\noindent\textbf{Figure 5: Confusion matrix showing classification accuracy}

\begin{figure}[htbp]
\centering
\includegraphics[width=0.9\columnwidth]{Figure5_Confusion_Matrix.png}
\label{fig:confusionmatrix}
\end{figure}

\noindent
Figure 5 shows the classification confusion table that is used to assess the classification performance. A high true positive and true negative value means that the various channel conditions and possible anomalies are correctly identified.
\FloatBarrier

\begin{table}[htbp]
\caption{Federated Learning vs Centralized Learning}
\centering
\begin{tabular}{|c|c|c|}
\hline
Parameter & Federated Learning & Centralized Learning \\
\hline
Data Storage & Local devices & Central server \\
\hline
Privacy & High & Low \\
\hline
Communication Cost & Moderate & High \\
\hline
Scalability & High & Limited \\
\hline
Security Risk & Lower & Higher \\
\hline
\end{tabular}
\end{table}

\section{Conclusion}

This paper has presented a federated learning-based system that improves the performance and security of quantum key distribution systems. The suggested solution combines Quantum Key Distribution and distributed learning methods allowing making smart predictions of the Quantum Bit Error Rate (QBER) and maintaining privacy of the data at all nodes. The framework enables the local training of models as well as sharing only updated models rather than raw data, which reduces communication overhead and enhances the scaling of distributed quantum communication settings. 
The results of the simulation indicate that the federated learning model is capable of the accurate prediction of QBER in varying channel conditions, thus enhancing the quality of the secure generation of keys. The proposed framework has similar accuracy in prediction and a higher level of privacy protection and less cost of communication as compared to the traditional centralized learning methods. The suggested system will be appropriate in the next-generation secure quantum communication networks due to these advantages. 
Further efforts will be devoted to the enhancement of the suggested framework by adding more sophisticated machine learning models and refining the process of aggregation in order to achieve a faster speed of convergence. The system can also be tested with more realistic quantum channel parameters (such as noise, attenuation and dynamic channel variations) in order to more fully test its functionality in realistic quantum communication circumstances. 
\FloatBarrier

\section*{Acknowledgment}

The authors thank the Department of Electronics and Communication Engineering, NIT Kurukshetra for research support.

\begin{thebibliography}{00}

\bibitem{b8}
L. Yuan, X. Chen, and T. Li, “Decentralized Federated Learning: A Survey and Perspective,” IEEE Internet of Things Journal, 2024.

\bibitem{b9}
F. Lin, C. Lee, and K. Chen, “Semi-Decentralized Federated Learning With Cooperative D2D Local Model Aggregations,” IEEE Journal on Selected Areas in Communications, 2021.

\bibitem{b10}
A. Mathur, S. Gupta, and M. Singh, “When Federated Learning Meets Quantum Computing: Survey and Research Opportunities,” IEEE Communications Surveys \& Tutorials, 2026.

\bibitem{b11}
 N. Huang, M. Dai, Y. Wu, T. Q. S. Quek, and X. Shen, 
“Wireless Federated Learning With Hybrid Local and Centralized Training: A Latency Minimization Design,” 
IEEE, Year.

\bibitem{b12}
 M. Chen, H. V. Poor, W. Saad, and S. Cui, 
“Convergence Time Optimization for Federated Learning Over Wireless Networks,” 
IEEE Trans. Wireless Commun., vol. 20, no. 4, pp. 2457–2471, Apr. 2021.

\bibitem{b13}
 E. Shalabi, 
“A Survey of Federated Learning Privacy Preservation Techniques for Malicious Behavior Detection,” 
J. Inf. Syst. Eng. Manag., vol. 10, no. 2, pp. XX–XX, 2025.

\bibitem{b14}
Nouhaila Innan, Muhammad Al-Zafar Khan, Alberto Marchisio, Muhammad 
Shafique, and Mohamed Bennai, “FedQNN: Federated Learning Using 
Quantum Neural Networks,” in Proceedings of the IEEE International 
Conference, IEEE, 2024.

\bibitem{b15}
D. K. Shareef, T. Shakeel, E. Hariprasad, A. Naresh, 
G. Naga Rama Devi, and N. V. Gangabathula, 
“Integrating Federated Learning in Quantum Computing for 
Secure and Scalable Model Training,” IEEE, 2024.
\bibitem{b16}
Lalitha K. P., Achanta Sree Rama Shriya, R. Kishore, and Krishna Lakshmi S., “Evaluating BB84 Protocol for Next-Generation Quantum Key Distribution Systems,” in Proceedings of an IEEE Conference, IEEE, 2024.

\bibitem{b17}
Gustav Nilsson, Martin Boldt, and Sadi Alawadi, “The Role of the Data Quality on Model Efficiency: An Exploratory Study on Centralised and Federated Learning,” in Proceedings of an IEEE Conference, IEEE, 2023.

\end{thebibliography}

\end{document}